% Chapitre 1 : Introduction
\chapter{Introduction}
\label{chap:introduction}

\section{Contexte Général}

L'analyse forensique des scènes de crime représente un domaine critique où la précision et la rapidité d'intervention peuvent avoir des conséquences significatives sur le cours d'une enquête. Traditionnellement, cette analyse repose sur l'expertise humaine, ce qui peut être sujet à des erreurs, des biais ou des limitations temporelles.

L'avènement de l'intelligence artificielle et du Deep Learning offre de nouvelles perspectives pour augmenter les capacités des experts forensiques. Les techniques de vision par ordinateur permettent désormais de détecter automatiquement des objets d'intérêt, d'extraire des informations textuelles et de générer des hypothèses préliminaires.

Notre projet, \textbf{AI Forensic Avatar}, s'inscrit dans cette démarche d'innovation en proposant un système complet d'analyse forensique assistée par l'IA, doté d'une interface utilisateur immersive basée sur un avatar 3D interactif.

\section{Problématique}

L'analyse de scènes de crime présente plusieurs défis majeurs :

\begin{enumerate}
    \item \textbf{Volume de données} : Les enquêtes modernes génèrent un nombre considérable d'images et de preuves visuelles qui nécessitent un traitement rapide.
    
    \item \textbf{Expertise limitée} : Les experts forensiques qualifiés sont rares et leur temps est précieux.
    
    \item \textbf{Subjectivité} : L'interprétation humaine peut varier d'un expert à l'autre, introduisant des biais potentiels.
    
    \item \textbf{Accessibilité} : Les outils d'analyse professionnels sont souvent coûteux et complexes à utiliser.
    
    \item \textbf{Communication} : La présentation des résultats d'analyse peut être technique et difficile à comprendre pour les non-spécialistes.
\end{enumerate}

\section{Objectifs du Projet}

Notre projet vise à développer une solution logicielle complète répondant aux objectifs suivants :

\subsection{Objectifs Techniques}

\begin{itemize}
    \item Implémenter un pipeline d'analyse d'images utilisant des modèles de Deep Learning état de l'art
    \item Développer une API REST robuste et performante avec FastAPI
    \item Créer une interface utilisateur moderne et intuitive avec React et Three.js
    \item Intégrer un système de streaming en temps réel pour les réponses de l'IA
    \item Mettre en place un avatar 3D avec synchronisation labiale pour une expérience immersive
\end{itemize}

\subsection{Objectifs Fonctionnels}

\begin{itemize}
    \item Permettre le téléchargement et l'analyse automatique d'images de scènes
    \item Détecter et classifier les objets présents sur les images
    \item Extraire le texte visible (plaques d'immatriculation, panneaux, documents)
    \item Générer des hypothèses forensiques basées sur les preuves détectées
    \item Offrir une interface conversationnelle avec l'avatar pour discuter des analyses
\end{itemize}

\section{Technologies Utilisées}

Le tableau \ref{tab:technologies} présente les principales technologies utilisées dans ce projet :

\begin{table}[H]
\centering

\label{tab:technologies}
\begin{tabular}{@{}llp{9cm}@{}}
\toprule
\textbf{Catégorie} & \textbf{Technologie} & \textbf{Rôle} \\
\midrule
\multirow{4}{*}{Backend} 
    & FastAPI & Framework web Python asynchrone \\
    & PostgreSQL & Base de données relationnelle \\
    & MinIO & Stockage objet compatible S3 \\
    & SQLAlchemy & ORM pour Python \\
\midrule
\multirow{4}{*}{IA/ML}
    & YOLOv8 & Détection d'objets en temps réel \\
    & Tesseract & Reconnaissance optique de caractères \\
    & Ollama/OpenAI & Génération de texte (LLM) \\
    & Groq PlayAI & Synthèse vocale (TTS) \\
\midrule
\multirow{4}{*}{Frontend}
    & React 19 & Bibliothèque UI JavaScript \\
    & Three.js & Rendu 3D WebGL \\
    & React Three Fiber & Intégration Three.js avec React \\
    & Zustand & Gestion d'état \\
\midrule
\multirow{2}{*}{DevOps}
    & Docker & Conteneurisation \\
    & Docker Compose & Orchestration multi-conteneurs \\
\bottomrule
\end{tabular}
\caption{Stack Technologique du Projet}
\end{table}

\section{Organisation du Rapport}

Ce rapport est organisé comme suit :

\begin{description}
    \item[Chapitre 2 - État de l'Art] Présente les concepts théoriques et les technologies existantes dans les domaines de la vision par ordinateur, du traitement du langage naturel et des interfaces 3D.
    
    \item[Chapitre 3 - Architecture Système] Décrit l'architecture globale du système, les choix de conception et les patterns utilisés.
    
    \item[Chapitre 4 - Implémentation Backend] Détaille l'implémentation du serveur FastAPI, les services IA et la gestion des données.
    
    \item[Chapitre 5 - Implémentation Frontend] Présente le développement de l'interface utilisateur React et l'intégration de l'avatar 3D.
    
    \item[Chapitre 8 - Conclusion] Résume les contributions et propose des perspectives d'amélioration.
\end{description}
